\section{Fundamentos de finanças: do preço ao retorno}
\label{sec:financas}

\subsection{Preço não é a melhor variável}
O dado bruto de uma ação é o \textbf{preço de fechamento} de cada dia. Mas o
preço, sozinho, é um número difícil de comparar: uma ação que custa R\$50 e
sobe R\$1 variou muito menos, proporcionalmente, do que uma de R\$5 que sobe
R\$1. Além disso, preços tendem a crescer ao longo dos anos (inflação,
crescimento), o que confunde um modelo que busca padrões estáveis.

\subsection{Retorno: a variação relativa}
\begin{definicao}[Retorno simples]
O retorno simples do dia $t$ é a variação percentual do preço:
\[
R_t = \frac{P_t - P_{t-1}}{P_{t-1}} = \frac{P_t}{P_{t-1}} - 1 ,
\]
onde $P_t$ é o preço de fechamento no dia $t$.
\end{definicao}

\subsection{Log-retorno: a versão que usamos}
Na prática trabalhamos com o \textbf{log-retorno}, o logaritmo natural da razão
entre preços consecutivos:
\[
r_t = \ln\!\left(\frac{P_t}{P_{t-1}}\right).
\]

\begin{intuicao}
Para variações pequenas (o caso quase sempre), $r_t \approx R_t$: uma alta de
1\% dá $r_t \approx 0{,}00995$. A vantagem do logaritmo é matemática e prática:
\begin{itemize}[noitemsep]
  \item \textbf{Somam-se no tempo.} O log-retorno de uma semana é a soma dos
        log-retornos diários. Com retornos simples teríamos que multiplicar
        $(1+R_1)(1+R_2)\dots$, bem menos cômodo.
  \item \textbf{Simetria.} Cair de 100 para 50 e voltar de 50 para 100 dá
        log-retornos $-0{,}69$ e $+0{,}69$ (opostos exatos); com retorno simples
        seriam $-50\%$ e $+100\%$, assimétricos.
  \item \textbf{Imune a fatores multiplicativos} como desdobramentos de ações
        (\emph{splits}), porque o log transforma multiplicação em soma.
\end{itemize}
\end{intuicao}

\begin{exemplo}
Preço passa de $P_{t-1}=20{,}00$ para $P_t=20{,}40$.
Retorno simples: $20{,}40/20{,}00 - 1 = 0{,}02$ (2\%).
Log-retorno: $\ln(20{,}40/20{,}00) = \ln(1{,}02) \approx 0{,}0198$. Quase iguais.
\end{exemplo}

\begin{lstlisting}[caption={Log-retorno em uma linha com pandas/numpy.}]
import numpy as np
# precos: serie diaria de fechamento (P_t)
log_ret = np.log(precos / precos.shift(1)).dropna()
\end{lstlisting}

\subsection{Volatilidade: o ``nível de agitação''}
\begin{definicao}[Volatilidade]
Volatilidade é o quanto os retornos oscilam. Mede-se, em geral, pelo
\emph{desvio-padrão} dos retornos (Seção~\ref{sec:estatistica}): retornos muito
espalhados $\Rightarrow$ alta volatilidade; retornos quase constantes
$\Rightarrow$ baixa volatilidade.
\end{definicao}

Mercados alternam \textbf{regimes}: períodos calmos (baixa volatilidade) e
períodos turbulentos (alta volatilidade, tipicamente em crises). Esse conceito
será central: uma anomalia é, grosso modo, um comportamento incompatível com o
regime recente.

\begin{naprat}
Usamos o \textbf{log-retorno diário do preço de fechamento} como única variável
de entrada do modelo. A escolha está registrada como decisão de projeto: é
invariante a \emph{splits} e estaciona a série (remove a tendência de
crescimento do preço), facilitando o aprendizado do que é ``normal''. Uma
consequência importante (Seção~\ref{sec:vazamento}) é que o período de treino
2010--2019, embora chamado de ``normal'', contém crises reais (recessão de
2014--2016, Lava Jato): ``normal'' aqui é \emph{relativo ao que o modelo viu}.
\end{naprat}
